\documentclass[sigconf]{acmart}
\settopmatter{printacmref=false}
\renewcommand\footnotetextcopyrightpermission[1]{}
\pagestyle{plain}
\begin{document}

\title{The Database as the Trusted Computing Base for LLM Agents:\\
Experience from TGMS}

\author{Xiaofei Zhang}
\affiliation{\institution{University of Memphis}\city{Memphis}\state{TN}\country{USA}}
\email{xiaofei.zhang@memphis.edu}

\begin{abstract}
Placeholder: existing agent--database interfaces ask an unreliable model to
generate a query, interpret partial results, perform additional computation,
and state conclusions. We argue agent-native databases should instead expose
a small, typed, costed operator algebra and return evidence-carrying results
whose downstream claims can be mechanically verified.
\end{abstract}

\maketitle

\section{The problem and the position}
% p1: failure example; TCB thesis; why SQL/Cypher-gen and RAG are
% insufficient; the three principles; contributions.

\section{Agent-native database architecture}
% p2: trust-boundary figure; plan IR; contract-bearing operator
% abstraction; LLM role vs DBMS role.

\section{Evidence and belief semantics}
% p3: bi-temporal correction example; execution trace; claim contract;
% truncation-taint example; support tuple.

\section{The TGMS prototype and lessons}
% p4: operator categories (4-5, not 13); implementation footprint;
% three lessons from live traffic; why each generalizes.

\section{Preliminary evidence}
% p5: end-to-end results table (frozen campaign); ablation table;
% scaling figure; honest limitations.

\section{Research agenda}
% p6: agent-facing algebra design; cost-based rewriting; evidence-aware
% execution; verification under approximation; bi-temporal provenance;
% multi-agent concurrency.

\end{document}
